% chapters/cover.tex — Make LLM-basic カバー
% レトロテクニカル（ブループリント/回路図スタイル）モノクロイラスト
% 1ページに正確に収まるように設計

\begin{titlepage}
\thispagestyle{empty}
\setlength{\parindent}{0pt}

% ============背景+フレーム(絶対位置)============
\begin{tikzpicture}[remember picture, overlay]
  % 背景（クリーム色の紙の感じ）
  \fill[paperbg] (current page.south west) rectangle (current page.north east);

  % 外枠（二重線）
  \draw[inkblue, line width=1.8pt]
    ($(current page.north west)+(1.2cm,-1.2cm)$)
    rectangle
    ($(current page.south east)+(-1.2cm,1.2cm)$);
  \draw[inkblue, line width=0.5pt]
    ($(current page.north west)+(1.4cm,-1.4cm)$)
    rectangle
    ($(current page.south east)+(-1.4cm,1.4cm)$);

  % 上：回路図パターン（ニューラルネットワークノード+接続線）
  \begin{scope}[shift={($(current page.north)+(0,-3.5)$)}, inkblue, line width=0.5pt]
    \foreach \x in {-4.5, -1.5, 1.5, 4.5} {
      \filldraw[inkblue] (\x, 0) circle (3pt);
    }
    \foreach \x in {-5, -2.5, 0, 2.5, 5} {
      \filldraw[inkblue] (\x, -1.2) circle (3pt);
    }
    \foreach \i in {-4.5, -1.5, 1.5, 4.5} {
      \foreach \j in {-5, -2.5, 0, 2.5, 5} {
        \draw[inkblue!25, line width=0.2pt] (\i, 0) -- (\j, -1.2);
      }
    }
    \filldraw[inkblue] (0, -2.4) circle (4pt);
    \draw[inkblue!60] (0, -1.2) -- (0, -2.4);
  \end{scope}

  % 下部：変圧器ブロックミニ図
  \begin{scope}[shift={($(current page.south)+(0,4.5)$)}, inkblue, line width=0.5pt]
    \draw[inkblue, fill=inkblue!5, rounded corners=2pt] (-3, 0) rectangle (3, 0.7);
    \node[inkblue, font=\small\sffamily\bfseries] at (0, 0.35) {Transformer Block};
    \draw[inkblue, fill=inkblue!10, rounded corners=2pt] (-2.4, 0.9) rectangle (2.4, 1.4);
    \node[inkblue, font=\scriptsize\sffamily] at (0, 1.15) {Multi-Head Attention};
    \draw[inkblue, fill=inkblue!10, rounded corners=2pt] (-2.4, -0.7) rectangle (2.4, -0.2);
    \node[inkblue, font=\scriptsize\sffamily] at (0, -0.45) {Feed-Forward};
    \draw[-{Stealth[length=4pt]}, inkblue!70] (0, 0.7) -- (0, 0.9);
    \draw[-{Stealth[length=4pt]}, inkblue!70] (0, -0.2) -- (0, 0);
  \end{scope}

  % 左：測定定規（レトロテクニカル感）
  \begin{scope}[shift={($(current page.west)+(1.8,2)$)}, inkblue!50, line width=0.3pt]
    \foreach \y in {0, 0.5, 1, 1.5, ..., 18} {
      \draw ($(0, -\y)$) -- ($(0.3, -\y)$);
    }
    \draw[inkblue, line width=0.5pt] (0, 0) -- (0, -18);
  \end{scope}
\end{tikzpicture}

% ============テキストコンテンツ（上から流れ）============
\vspace*{0.3\textheight}

\begin{center}
% メインタイトル
{\fontsize{56pt}{68pt}\selectfont\bfseries\sffamily\color{inkblue} Make LLM}\\[0.3cm]

% サブタイトル
{\fontsize{32pt}{38pt}\selectfont\bfseries\sffamily\color{inkgray} basic}\\[1.0cm]

% 区切り線
\begin{tikzpicture}
  \draw[inkblue, line width=0.8pt] (0,0) -- (9,0);
\end{tikzpicture}\\[0.4cm]

% 説明フレーズ
{\large\itshape\sffamily\color{inkgray}トルクナイザーから変圧器まで\\[0.2cm]
最初から自分で作る大型言語モデル}\\[0.5cm]

\begin{tikzpicture}
  \draw[inkblue, line width=0.8pt] (0,0) -- (9,0);
\end{tikzpicture}\\[0.6cm]

% 著者/出版社（表紙の一番下）
{\large\sffamily\bfseries\color{inkblue} dirmich}\\[0.2cm]
{\normalsize\sffamily\color{inkgray} Highmaru Press}
\end{center}

\end{titlepage}

% 表紙 次ページへ強制移動(ミックス防止)
\clearpage
